\chapter{Adaptive Liquidity Provision via Reinforcement Learning}
\label{ch:rl_strategy}

\section{Introduction}
The inherent non-stationarity of financial markets, characterized by shifting volatility regimes and non-linear microstructural dependencies, necessitates a quoting strategy that possesses high adaptive capacity. This chapter details the formulation of the reinforcement learning (RL) framework used to optimize the market-making policies defined in the stochastic control context of Chapter \ref{ch:microstructure}. We focus on the transformation of the continuous quoting problem into a discrete-time Markov Decision Process (MDP) and delineate the learning algorithm employed to derive robust execution policies.

\section{Markov Decision Process (MDP) Formulation}
We model the market-making task as a discrete-time MDP defined by the tuple $(S, A, P, R, \gamma)$.

\subsection{State Space Representation ($S$)}
The state vector $s_t \in S$ is engineered to capture a multidimensional summary of the instantaneous Limit Order Book (LOB) state and the market maker's internal constraints. The features include:
\begin{equation}
    s_t = \left[ \rho_t, \sigma_t, q_t, \psi_{t-1}, \Delta S_t, \text{imb}_{vol} \right]
\end{equation}
where:
\begin{itemize}
    \item $\rho_t$: Normalized Order Book Imbalance.
    \item $\sigma_t$: Rolling volatility estimate (GARCH-like feature).
    \item $q_t$: Current net inventory position.
    \item $\psi_{t-1}$: The spread utilized in the previous time step.
    \item $\Delta S_t$: Short-term mid-price innovation.
\end{itemize}

\subsection{Action Space and Quoting Policy ($A$)}
The agent's action $a_t \in A$ corresponds to the selection of bid and ask offsets $(\delta^b, \delta^a)$ relative to the reservation price $r_t$ derived in Equation 3.10. To maintain computational efficiency within the FPGA fabric, we discretize the action space into a finite set of spread-shift commands $(\Delta \delta)$. This discretization allows the model to "tilt" or "widen" the quoting envelope in response to detected regime shifts.

\subsection{Risk-Adjusted Reward Function ($R$)}
The objective of the RL agent is to maximize terminal utility by balancing the capture of the bid-ask spread against the costs of inventory holding and adverse selection. The reward function $r_t$ is formulated as:
\begin{equation}
    r_t = \Delta PnL_t - \eta \cdot \text{Var}(q_t) - \zeta \cdot \mathbf{I}(|q_t| > q_{max})
\end{equation}
where $\Delta PnL_t$ is the change in mark-to-market wealth, $\eta$ is the inventory risk aversion coefficient, and $\zeta$ is a heavy penalty for violating the hard risk limits enforced by the system's safety envelope.

\section{Learning Algorithm: Proximal Policy Optimization (PPO)}
To ensure stable policy updates in the noisy HFT environment, we utilize the **Proximal Policy Optimization (PPO)** algorithm \citep{Schulman2017}. PPO is particularly suited for financial applications due to its clipped surrogate objective, which prevents large, destructive updates to the policy parameters that could lead to catastrophic losses during the training phase.

The objective function minimized by the agent is:
\begin{equation}
    L^{CLIP}(\theta) = \hat{\mathbf{E}}_t \left[ \min(r_t(\theta) \hat{A}_t, \text{clip}(r_t(\theta), 1-\epsilon, 1+\epsilon) \hat{A}_t) \right]
\end{equation}
where $r_t(\theta)$ is the probability ratio between the new and old policies, and $\hat{A}_t$ is the estimated advantage.

\section{Hardware-Centric Model Optimization}
The transition from a high-level RL model to a hardware-executable core requires significant numerical optimization.

\subsection{Fixed-Point Quantization}
While RL models are typically trained using 32-bit floating-point precision, hardware execution on FPGAs is most efficient using fixed-point arithmetic. We employ **Quantization-Aware Training (QAT)** to ensure that the policy's performance does not degrade when converted to 16-bit fixed-point representations. This process allows the hardware DSP48 slices to perform the necessary Multiply-Accumulate (MAC) operations within a single clock cycle.

\subsection{Pruning and Sparsity}
To further minimize the inference latency, we apply iterative pruning to the neural network weights. By removing non-contributing synaptic connections, we reduce the total number of hardware gates required for implementation, facilitating the sub-15ns inference target discussed in Chapter \ref{ch:architecture}.

\section{Conclusion}
This chapter has detailed the formulation of the adaptive quoting strategy as an RL problem. By synthesizing stochastic control theory with the PPO algorithm and hardware-centric quantization, we create a policy that is both intelligent and execution-ready. The following chapters delineate the hardware-software architecture used to operationalize this policy at wire speed.
